% ============================================================================
% CHAPTER 2 - LITERATURE REVIEW
% File: docs/thesis/theory/ch2_grid_access_degradation.tex
% Grid access, external infrastructure costs and lifetime degradation.
% Operational values: docs/thesis/methods/group3_connection_degradation.tex
% ============================================================================

\section{Grid Access and External Infrastructure Costs}
\label{sec:grid_access_review}

The cost of a utility-scale solar or wind plant depends not only on its
resource and its equipment but also on the infrastructure needed to
deliver its output to the transmission system. Where good resources lie
far from the network, the distance to a suitable point of
interconnection (POI) introduces a cost that varies strongly in space.
Empirical analyses of interconnection requests in the United States show
that these costs differ widely between projects and that network
upgrades beyond the dedicated connection can account for a large share
of the total \cite{seel2023,miso2024}; transmission planning studies in
Australia reach the same conclusion for regional expansion options
\cite{aemo2023}. A single average connection cost applied to every site
therefore misrepresents the relative merit of locations. Spatial
supply-curve models such as NREL's reV instead combine resource, land
availability and the distance to transmission into site-specific
connection costs \cite{maclaurin2021}.

\subsection{Points of Interconnection and Open Grid Data}
\label{subsec:poi_review}

A POI is the interface between a plant and the shared network. For
plants of tens to hundreds of megawatts it is normally a transmission
substation able to accommodate the voltage, switching, protection and
power transfer involved, not an arbitrary point on a line. OpenStreetMap
(OSM) represents substations and lines with a nominal voltage tag, which
allows assets to be filtered by voltage level
\cite{osmPowerTag,osmVoltageTag}. Its completeness, however, varies
between countries. Arderne et al.\ combined OSM with night-time lights
and routing to predict the global power network, including medium-voltage
lines that are not mapped \cite{arderne2020}. Such predictive layers cover
more territory but do not record the voltage, status or capacity of each
segment, whereas tagged substations are fewer but interpretable.
Initiatives to improve open transmission data acknowledge this trade-off
\cite{oetMapYourGrid2025}. A transparent hierarchy follows: verified
high-voltage substations as the primary target, and predictive layers as
a sensitivity.

\subsection{From Distance to Route Length}
\label{subsec:route_review}

The straight-line distance to a POI is a first-order measure of
proximity. Real lines and roads deviate from it because of terrain,
protected areas, land tenure, existing infrastructure and right-of-way
requirements. Screening studies that do not route each connection
correct the straight-line distance with a tortuosity factor, the ratio of
route length to geometric distance; least-cost-path routing or measured
corridor lengths replace it when available
\cite{maclaurin2021,miso2024,aemo2023}. The same reasoning applies to the
access road from a site to the existing road network, which is distinct
from the internal roads already counted in plant construction.

\subsection{Cost Boundaries and Unit Costs}
\label{subsec:cost_boundary_review}

A plant's capital cost and its external connection must be separated
explicitly, or the same equipment is counted twice. NREL's Annual
Technology Baseline excludes grid connection from overnight capital cost
while including on-site electrical equipment such as the plant
switchyard; spur lines, the POI and nominal network upgrades are treated
separately \cite{nrelATB2024definitions}, as they are in reV
\cite{maclaurin2021}. The practical boundary places collection networks,
collector substations, step-up transformers and the switchyard inside the
plant, and the dedicated high-voltage spur line, the terminal bay at the
receiving substation and off-site access works outside it. The label
``total installed cost'' of a cost survey does not by itself establish
where its boundary lies, which must be read from the source.

Unit costs of lines and substations depend on voltage, circuit
configuration, conductor, structures, terrain and right of way.
Component-based guides such as MISO's give voltage-specific estimates
\cite{miso2024}, and AEMO's materials do the same for Australian
conditions \cite{aemo2023}. Their numbers do not transfer directly
between countries, where labour, materials, standards, taxes and
regulation differ, so regulatory benchmarks take precedence where they
exist: the indicative capital costs of the Indian Central Transmission
Utility for dedicated transmission lines \cite{ctuilIndicative2025}, the
reference price bank of the Brazilian regulator \cite{aneelBPR}, and the
unit investment cost indicators compiled by ACER from European
transmission projects \cite{acer2023uic}. A review of World Bank-financed
projects found unit costs well below those reported in the literature and
a wide spread within each voltage class, and warns that such figures give
only an order of magnitude \cite{esmap2026transmission}. A terminal bay
at an existing substation, the incremental switching and protection for
one feeder, is much cheaper than a new substation; broader reinforcements
such as transformers or new circuits are a separate cost.

\section{Lifetime Performance Degradation}
\label{sec:degradation_review}

Assuming constant output over a multi-decade life overstates discounted
energy and therefore understates the levelized cost. Degradation differs
by technology and location: PV loses performance through encapsulant
discoloration, delamination, corrosion, cell and interconnect damage,
potential-induced degradation and failures of balance-of-system
components, driven by temperature, humidity and ultraviolet exposure
\cite{jordan2017}; wind turbines lose performance through blade erosion,
aerodynamic changes, drivetrain wear and rising maintenance needs
\cite{hamilton2020,kim2025}.

For PV, reported rates vary with technology, mounting, climate and
measurement method. A recent meta-analysis of 610 observations from 80
studies found a median of about 0.94\,\%/yr and a mean of 1.1\,\%/yr
\cite{straubmuck2025}. Fleet data from the United States show the role of
temperature: Jordan et al.\ report a median system performance-loss rate
of 0.75\,\%/yr, about 0.48\,\%/yr in cooler and 0.88\,\%/yr in hotter
climates, and find high-efficiency technologies broadly comparable to
conventional silicon \cite{jordan2022}. Their climates are the
temperature zones of Karin et al., which bin the Arrhenius-weighted
equivalent temperature of an open-rack module (activation energy
1.1\,eV) in 5\,$^\circ$C intervals, from T1 below 14\,$^\circ$C to T10
\cite{karin2019}. The zones are a stress index derived from hourly
weather rather than a climate class, which makes them transferable to
other regions, but the rates themselves remain tied to the fleet in which
they were measured, whose hottest zone is T7. A low temperature coefficient of power
describes an instantaneous, reversible loss and is not evidence of slower
irreversible degradation.

For wind, Hamilton et al.\ analysed 917 plants in the United States and
found performance declines over the first ten years of about
0.17\,\%/yr for newer turbine vintages and 0.53\,\%/yr for older ones
\cite{hamilton2020}. They also observed a drop when plants lose
production tax-credit eligibility after ten years, an economic effect of
the United States policy rather than physical ageing. Kim and Kim report
about 0.72\,\%/yr for wind farms in complex terrain \cite{kim2025}. The
evidence supports vintage-specific rates for wind, with higher values
kept for demanding sites, and climate-specific rates for PV.
